\section{Relación entre sistemas dinámicos, topología y cálculo fraccionario}
\label{sec:programa-estudio-bibliografico}

El estudio de los atractores ocultos combina tres estructuras: la ecuación
que genera las trayectorias, el espacio donde evolucionan y los operadores
que incorporan su memoria. En orden entero, la unicidad transforma el PVI
en un flujo de puntos; en Caputo, Volterra permite levantarlo a una evolución
de perfiles. La tabla organiza las conexiones matemáticas que desarrollamos.

\begin{table}[H]
\centering
\small
\begin{tabularx}{\textwidth}{@{}L{0.10\textwidth}YYY@{}}
\toprule
 & Sistemas dinámicos & Topología y geometría & Cálculo fraccionario \\
\midrule
1 & Campo, solución, flujo y equilibrio & Espacio de fases, continuidad,
tangencia y curvas integrales & Gamma, Mittag--Leffler, Caputo y Volterra \\
2 & Linealización, estabilidad y bifurcación & Índice, grado, variedades
invariantes y transversalidad & Matignon, estado funcional y memoria \\
3 & Conjuntos límite, atracción, caos transitorio y predictibilidad &
Poincaré, Floquet, separatrices y crisis & Retorno funcional,
recurrencia y continuación con memoria \\
4 & Localización y tiempos de escape & Puntos perpetuos, superficies
críticas, borde de cuenca y aislamiento & Función descriptiva,
continuación causal, defecto de Volterra y convergencia \\
\bottomrule
\end{tabularx}
\end{table}

Las nociones de continuidad y compacidad intervienen tanto en la dependencia
respecto de las condiciones iniciales como en los argumentos de punto fijo.
La estabilidad espectral describe el entorno de un equilibrio; los mapas de
retorno y las cuencas incorporan trayectorias no estacionarias. En la
formulación fraccionaria mantenemos estas conexiones y sustituimos el
estado puntual por el perfil de memoria cuando la propiedad estudiada exige
una evolución autónoma completa.
